\section{Semantic Repair and Equation-Generated Geometry}\label{sec:4}

This section constructs the two complexes of SAGA independently. The
semantic side is generated from semantic atoms and a repair relation, the
equation side from an equation system; each produces its coefficient,
complex, and residual from its own primary data alone. Neither
construction refers to the other, and the comparison of
Section~\ref{sec:5} --- the existence of the map, its isomorphy, the
correspondence of residual classes --- is the conclusion of a theorem,
not a repetition of the constructions. The single map connecting the two,
$\chi^E$ (Proposition~\ref{prop:chiE}), is a construction of a comparison
input for Section~\ref{sec:5} and is used in the construction of neither
complex.

\subsection{The cover-relative \v{C}ech complex}\label{sec:41}

Let $F$ be an abelian-group-valued presheaf on $S_X$. Define the
three-term cochain complex relative to $\cU$ by
\[
C^0(\mathcal U,F)=\prod_iF(U_i),\qquad
C^1(\mathcal U,F)=\prod_{i<j}F(U_{ij}),\qquad
C^2(\mathcal U,F)=\prod_{i<j<k}F(U_{ijk})
\]
(the products run over nonempty intersections), with the differentials on
ordered indices
\[
(\delta^0a)_{ij}=a_j|_{U_{ij}}-a_i|_{U_{ij}},
\qquad
(\delta^1c)_{ijk}=c_{jk}|_{U_{ijk}}-c_{ik}|_{U_{ijk}}+c_{ij}|_{U_{ijk}}.
\]
A direct computation gives $\delta^1\delta^0=0$, so the cover-relative
$H^1$ is defined by
\[
\check H^1(\mathcal U,F)=\ker\delta^1/\operatorname{im}\delta^0.
\]

\subsection{The semantic-side construction}\label{sec:42}

\paragraph{Semantic repair presentation.}
To each context $V$ of $S_X$, assign the set $\Lambda(V)$ of semantic
atoms distinguished on $V$, a projection $\pi_V:\Lambda(V)\to\At$ giving
the underlying Atom of each semantic atom, and a subset
$S(V)\subseteq\Lambda(V)$ of \textbf{supported semantic atoms} available
for repair. For each context morphism $V'\to V$ there is a functorial
restriction map $\Lambda(V)\to\Lambda(V')$ that preserves support (if
$\lambda\in S(V)$ then $\lambda|_{V'}\in S(V')$), and the projection
preserves the underlying Atom:
\[
\pi_{V'}(\lambda|_{V'})=\pi_V(\lambda).
\]
Write $\Fsem(V)$ for the free abelian group on $S(V)$; its elements
(formal finite sums of supported atoms) are called \textbf{repair
words}. On each $V$, select a restriction-stable subgroup
\[
\Rrep(V)\subset \Fsem(V)
\]
(the \textbf{local repair relation}). This is all of the primary data of
the semantic side. The comparison core
(\S\ref{sec:53}--\S\ref{sec:55}) uses these data restricted to the cover
intersection diagram.

\paragraph{Semantic coefficient.}
The quotient
\[
\Msem(V)=\Fsem(V)/\Rrep(V)
\]
carries a restriction induced by restriction-stability, and
$V\mapsto \Msem(V)$ is an abelian-group-valued presheaf on $S_X$. The
complex of \S\ref{sec:41} with coefficients $\Msem$,
$C^\bullet_{\mathrm{sem}}(\cU):=C^\bullet(\cU,\Msem)$, is the semantic
complex.

\paragraph{Affine semantic repair system.}
Let $\Psem$ be a presheaf of semantic local repair states on $S_X$. On
each $V$, $\Fsem(V)$ acts on $\Psem(V)$, restriction commutes with the
action, and on each $V$ of the cover intersection diagram the following
three conditions hold:
\begin{itemize}
\item \textbf{action soundness}: words in $\Rrep(V)$ act as the
identity;
\item \textbf{stabilizer completeness}: if $\Psem(V)$ is nonempty, every
word acting as the identity belongs to $\Rrep(V)$;
\item \textbf{local transitivity}: any two states are carried to each
other by the action of some word.
\end{itemize}
By soundness the action descends to an action of $\Msem(V)$; by
completeness this action is free; by transitivity it is transitive on
nonempty $\Psem(V)$. Therefore, on each $V$ of the cover intersection
diagram, a nonempty $\Psem(V)$ is an $\Msem(V)$-torsor (an affine
space). A torsor is a space without a distinguished origin in which,
instead, the difference of any two states is measured uniquely as an
element of the coefficient group. The construction of residuals below
uses only this difference, never the states themselves.

\paragraph{Semantic residual.}
Choosing a selected local repair atlas $\{p_i\}$, the torsor difference
on each nonempty overlap determines a unique
\[
r_{\mathrm{sem},ij}=p_j|_{U_{ij}}-p_i|_{U_{ij}}\in \Msem(U_{ij}),
\]
giving the \textbf{semantic residual} $\rsem\in C^1_{\mathrm{sem}}(\cU)$.
The residual $\rsem$ is a cocycle: on $U_{ijk}$ we have
$p_k=r_{jk}+p_j=r_{jk}+r_{ij}+p_i$ and $p_k=r_{ik}+p_i$, and freeness of
the action gives $r_{jk}-r_{ik}+r_{ij}=0$.

The class $[\rsem]\in H^1_{\mathrm{sem}}(\cU)$ does not depend on the
choice of atlas (\textbf{choice independence}): for another atlas
$p'_i=a_i+p_i$ the same torsor computation gives
$r'_{\mathrm{sem}}=\rsem+\delta^0a$. Consequently, the vanishing
$[\rsem]=0$ is equivalent to the existence of a correction
$a\in C^0_{\mathrm{sem}}(\cU)$ such that the corrected atlas
$(-a_i)+p_i$ agrees on all pairwise overlaps (the existence of a
\textbf{matching correction}). This paraphrase provides the intuition
for coboundaries: a coboundary $\delta^0a$ is the difference produced by
merely changing the reference (the atlas) within each chart. A zero
class is an inconsistency that can be removed by per-chart changes of
reference; a nonzero class is an inconsistency that no combination of
local reference changes removes, and that remains on a closed loop.

\subsection{The equation-side construction}\label{sec:43}

The equation side constructs the \textbf{geometric \v{C}ech complex}
$C^\bullet_E(\cU):=C^\bullet(\cU,\QE)$ with coefficients in the $\QE$ of
Theorem~\ref{thm:quotient}. The \textbf{equation-lift system} $\PE$ is a
presheaf on $S_X$ such that on each intersection $V$, the nonempty
$\PE(V)$ carries a free and transitive action of $\QE(V)$, and
restriction commutes with the action. From a selected local lift atlas
$\{e_i\}$, the difference on overlaps generates the geometric-side
residual
\[
r_{E,ij}=e_j|_{U_{ij}}-e_i|_{U_{ij}}.
\]
By the same torsor argument as in \S\ref{sec:42}, $\rE$ is a cocycle and
its class does not depend on the choice of atlas.

\begin{proposition}[equation semantic realization: construction of $\chi^E$]\label{prop:chiE}
To each supported semantic atom $\lambda\in S(V)$ on each cover
intersection $V$, assign a required equation index $i_\lambda$ and a
local architecture reading $A_\lambda$, compatibly with restriction
(selected; for a face restriction $V'\to V$ of the cover intersection
diagram, $i_{\lambda|_{V'}}=i_\lambda$ and
$A_{\lambda|_{V'}}=A_\lambda$). Then
\[
\chi^E_V(\lambda):=[\epsilon_{V,A_\lambda,i_\lambda,\pi_V(\lambda)}]\in \QE(V)
\]
is restriction-natural with respect to face restrictions, i.e.
\[
\chi^E_{V'}(\lambda|_{V'})=\chi^E_V(\lambda)|_{V'}.
\]
\end{proposition}

\begin{proof}
The indices $i_\lambda$ and readings $A_\lambda$ are chosen compatibly
with restriction, and $\pi$ preserves the underlying Atom
(\S\ref{sec:42}). Hence the residual at $\lambda|_{V'}$ is the
$\epsilon$ evaluated on $V'$ with the same reading, index, and support
Atom, and by the residual restriction naturality of
Theorem~\ref{thm:quotient} its class coincides with the restriction of
$\chi^E_V(\lambda)$.
\end{proof}

$\chi^E$ is the canonical instance, constructed from the equation system
and displayed readings, of the restriction-natural map that assumption~3
of Theorem~\ref{thm:central} requires; it shows that the starting point
of the comparison map is generated from nothing more than ``which
failure, of which reading of which equation''.
